\section{Partial results at the next heterotic tree order}\label{sec:nexttree}

The amplitude in Theorem~\ref{thm:amplitude} includes the complete
Yang--Mills tower. A further coefficient is required to include the
other heterotic interactions. We identify its perturbative order, bound
the quadratic potential mixing with K3 shape deformations, and calculate
the universal massless spin-two contribution. The remaining
coefficient is not evaluated here.

\subsection{Order counting and the required terms}
Remove the common gravitational prefactor and write the action in
reference internal units as
\begin{equation}\label{eq:treeorders}
 S[Y,\mathcal A]=S_0[Y]+\eta S_1[Y,\mathcal A]
             +\eta^2 S_2[Y,\mathcal A]+\Oformal(\eta^3),
\end{equation}
where $Y$ denotes the metric, dilaton and two-form. The term $S_1$
contains both Yang--Mills and tangent curvature squared. In the field
scheme of \eqref{eq:heterotic-action}, $S_2$ includes the terms generated
by the recursive Chern--Simons definition of $H$ and the composite
Hull connection. The absence of a new independent bosonic invariant
at this order does not set $S_2$ to zero
\cite[section~3.1]{aqs}\cite[section~2.1]{mcorist_yin,mcorist_picard}.

Each canonical gauge or connection-scalar external leg removes a factor
$\eta^{1/2}$ in these units. A bare term $\eta^nS_n$ with four such
legs has order $\eta^{n-2}$. The leading Yang--Mills amplitude has
order $\eta^{-1}$; its first relative-$\eta$ correction requires the
complete bare action through $\eta^2$. The same count follows for
exchange diagrams after canonical normalization of internal fields.
Independent bare $\eta^3$ terms affect a later coefficient at fixed
deformation and momentum. This order count gives no estimate uniform
in a vanishing Higgs mass or near a pole.

One can organize the calculation by the stationary action. Let $Y_0$
be the leading product background, $K$ the gauge-fixed Hessian of $S_0$,
and $J[\mathcal A]=\delta_Y S_1[Y_0,\mathcal A]$. With constraints,
zero modes and propagator boundary conditions specified, eliminating
the neutral fluctuations to this order would give
\begin{equation}\label{eq:stationarycorrection}
 W[\mathcal A]=S_2[Y_0,\mathcal A]
       -\tfrac12\langle J[\mathcal A],K_F^{-1}J[\mathcal A]\rangle.
\end{equation}
This follows by inserting the stationary displacement
$-\eta K_F^{-1}J$ into the quadratic leading action and its source.
The massless poles in $K_F^{-1}$ are retained; the expression is not
a local elimination of massless fields. At zero external momentum,
the compatible complementary inverse is used, with background zero
modes fixed by the stated normalization.

Formula~\eqref{eq:stationarycorrection} is an organization of the
missing calculation, not a computed amplitude. If $J_d$ is the
degree-$d$ fluctuation part, its quartic term involves
$J_0K_F^{-1}J_4$, $J_1K_F^{-1}J_3$ and
$\tfrac12J_2K_F^{-1}J_2$. The first two include the corrected background
and linear mixing. One must also insert the cubic part of $W$ into
Yang--Mills exchange diagrams, insert its quadratic part into internal
propagators, and correct the external modes and residues. Adding only
neutral exchange between two scalar-pair sources omits these terms.
The scalar constraints in section~\ref{sec:scalar} and the vector pencil
in section~\ref{sec:coupled} supply parts of this reduction, but neither
determines the complete cubic and quartic vertices.

\subsection{Quadratic mixing with K3 metric deformations}
One part of the coupled reduction admits a relative estimate from the
same harmonic energy bound. Consider a volume-preserving infinitesimal
hyperk\"ahler metric deformation in harmonic gauge, modulo
diffeomorphisms and rotations of the parallel self-dual triple.
Its period variations are harmonic ASD two-forms
$\beta_\alpha=(\partial_\tau\omega_\alpha)^-$; this is the usual
$\mathbb R^3\otimes\mathcal H^-_2(S)$ description of the 57-dimensional
space of K3 shape tangents \cite{gik}. Normalize
$\sum_\alpha\|\beta_\alpha\|^2\le1$.

At fixed connection, put $f_\alpha=\partial_\tau\mu_\alpha$ and
$c_{n\alpha}=\langle p_n,f_\alpha\rangle$. Up to the fixed component
signs in \eqref{eq:infinitemoments}, the source is
$f_\alpha=*_0(F\wedge\beta_\alpha)$. The variation of the outer
Hodge star contributes zero because $F\wedge\omega_\alpha=0$ at the
ASD background, and the self-dual part of the period variation wedges
to zero with $F$. Since $\beta_\alpha$ is harmonic and ASD,
\begin{equation}\label{eq:shapesource}
 |c_{n\alpha}|
 =\left|\int_S\BF(p_n,F)\wedge\beta_\alpha\right|
 =\left|\langle\HH_2\BF(p_n,F),\beta_\alpha\rangle\right|.
\end{equation}

\begin{lemma}\label{lem:shapemixing}
For each low positive mode,
\begin{equation}\label{eq:shapemixing}
 \sum_\alpha|c_{n\alpha}|^2\le Cr^2\mu_n.
\end{equation}
The mixed Hessian of $\tfrac12\sum_\alpha\|\mu_\alpha\|^2$
between the unit scalar normal $n_{n\alpha}$ and this metric tangent
is $\sqrt{\mu_n}c_{n\alpha}$, hence is $O(r\mu_n)$.
The normal gauge-field adjustment to the linearized ASD equation is
\begin{equation}\label{eq:shapelift}
 \dot{\mathcal A}_\perp
 =-\sum_{n>0,\alpha}\frac{c_{n\alpha}}{\sqrt{\mu_n}}n_{n\alpha}.
\end{equation}
Its low component is $O(r)$ and its complementary component is $O(1)$
in every fixed spatial Sobolev norm, uniformly in $b/a$.
\end{lemma}
\begin{proof}
Equations~\eqref{eq:shapesource} and~\eqref{eq:relativeharmonic}
give \eqref{eq:shapemixing} by Cauchy--Schwarz. They also show that
the source has zero projection on every parallel zero-form.
At zero moment level the potential Hessian is the pairing of first
moment derivatives. Equation~\eqref{eq:quaternionicgram} gives the
displayed mixed entry and verifies that \eqref{eq:shapelift} solves
$D\mu_\alpha(\dot{\mathcal A}_\perp)+f_\alpha=0$.

For the full adjoint ASD complex the moment-map cokernel consists of
three copies of the parallel algebra; the preceding vanishing removes
each compatibility obstruction. The full lift includes the spectator
$B$ connection's curvature source and its gapped response. In a
trace-free determinant-fixed stable $E$ block the obstruction space
is zero; restricting the lift to that block means restricting the
source as well.

Orthonormality of the normals gives
\[
 \|\dot{\mathcal A}_{\perp,\mathrm{low}}\|^2
 =\sum_{n\,\mathrm{low},\alpha}\frac{|c_{n\alpha}|^2}{\mu_n}
 \le Cr^2.
\]
Each unit low normal solves the uniformly elliptic ASD one-form
square with bounded eigenvalue. Its coefficients are uniformly
smooth in the fixed unitary presentation, so elliptic regularity
bounds the normal in every fixed $H^s$ norm. The finite low sum
therefore has the same $O(r)$ bound in $H^s$; dividing a spatial
estimate for $Dp_n$ by a small $\sqrt{\mu_n}$ is unnecessary.
For the complementary part, Parseval and the fixed positive gap
give a uniform $L^2$ bound. The curvature and normalized harmonic
$\beta_\alpha$ have uniformly bounded spatial derivatives, so the
source does also. The complementary ASD equation and elliptic
regularity then give its fixed-$H^s$ bounds.
\end{proof}

In the common normalization of \eqref{eq:treeorders}, a canonical
gauge normal carries $\eta^{-1/2}$ and a geometry mode has leading
order-one normalization. Thus this potential block has a low mixed
entry $O(\sqrt\eta\,r\mu_n)$ and a static canonical normal
admixture $O(\sqrt\eta\,r)$. This controls one quadratic contribution.
It does not bound parameter derivatives of \eqref{eq:shapelift},
the full corrected kinetic form, or a two-external-scalar neutral
source at nonzero momentum. Those are still needed for the complete
four-point coefficient.

\subsection{The massless graviton and the angular domain}
Use the six-dimensional Einstein frame normalized to agree with the
product-frame external metric at the fixed leading vacuum. Define
\begin{equation}\label{eq:Einsteinnormalization}
 S_{\rm EH}=\frac1{2\kappa_6^2}\int\sqrt{-g}\,R,
 \qquad g_{\mu\nu}=\eta_{\mu\nu}+2\kappa_6h_{\mu\nu}.
\end{equation}
The curvature convention in this Lorentzian action gives the usual
healthy graviton propagator. For canonical orthogonal massless scalars,
the minimally coupled graviton vertex is diagonal in species. Hence
$\phi\chi\to\phi\chi$ has only $t$-channel graviton exchange at
this order. It is
\begin{equation}\label{eq:gravitonamplitude}
 \mathcal M_{\rm grav}=\kappa_6^2\frac{su}{t}.
\end{equation}
For completeness, the conserved scalar stress at a vertex is
$T_{\mu\nu}(p,p')=p_\mu p'_\nu+p_\nu p'_\mu-\eta_{\mu\nu}p\cdot p'$.
The conserved-source propagator numerator in $D$ dimensions is
\[
 P_{\mu\nu,\rho\sigma}
 =\tfrac12(\eta_{\mu\rho}\eta_{\nu\sigma}
           +\eta_{\mu\sigma}\eta_{\nu\rho})
       -\frac{\eta_{\mu\nu}\eta_{\rho\sigma}}{D-2}.
\]
On-shell contractions give
$T_1PT_2=(s^2+u^2-t^2)/2=-su$. In mostly-plus signature the exchanged
momentum has $q^2=-t$, and the two vertices and propagator give
$\mathcal M=\kappa_6^2T_1PT_2/q^2=-\kappa_6^2T_1PT_2/t$, proving
\eqref{eq:gravitonamplitude} also in $D=6$. Conservation removes
longitudinal gauge terms.

Restoring the common normalization of \eqref{eq:heterotic-action},
the gauge term is
\[
 -\frac{\alpha'e^{-2\Phi_0}}{8\kappa_{10}^2}
                    \BF(F_{MN},F^{MN}).
\]
Comparison with \eqref{eq:YMaction} and integration over the leading
volume $V_0L^4$ give
\begin{equation}\label{eq:gravitycoupling}
 g_{10}^2=\frac{2\kappa_{10}^2e^{2\Phi_0}}{\alpha'},\qquad
 \kappa_6^2=\frac{\kappa_{10}^2e^{2\Phi_0}}{V_0L^4}
           =\frac34\alpha'e^2.
\end{equation}
These are the leading background normalizations needed for this
coefficient. If $|t|\ge\delta E^2$, then
\begin{equation}\label{eq:gravitybound}
 |\mathcal M_{\rm grav}|
 \le\frac{3}{4\delta}e^2\eta(EL)^2.
\end{equation}
At $E\asymp b/L$ this is $O(e^2\eta b^2)$, uniformly in $b/a$.
It is an absolute bound relative to the natural amplitude scale $e^2$,
not a relative bound against a Yang--Mills amplitude that may vanish.

The massless pole requires a change to the domain of any full-amplitude
comparison. For sufficiently small nonzero $s$, the domain
\eqref{eq:amplitudedomain} allows $t\to0$ with $u\to-s$.
The graviton term then has the nonzero residue $-\kappa_6^2s^2$.
Thus the pure-Yang--Mills comparison cannot be extended to the
unsubtracted full amplitude on precisely the same forward-allowed
domain by appending a bounded error. Angular separation from massless
poles, or a specified gauge-invariant pole subtraction, is necessary.
This restriction does not by itself control the other terms of
\eqref{eq:stationarycorrection}.

The unresolved part is the complete constrained neutral-field source,
its recursively induced contacts, the corrected scalar external states,
and their combined exchanges and residues. Uniform estimates would
also have to control small gauge inverses and mass insertions in this
sum. They do not follow from the fixed complementary gaps alone.
Even a complete calculation of this coefficient would remain distinct
from a bound on the full finite-$\alpha'$ remainder.
